% page 635 of the facsimile (image p-634 of the scan)
% block 167.6 x 254.9 mm ; left margin 34.4 mm
\pgc{12.700mm}{0.000mm}{
\l{}
\l{}
\l{}
\l{\cel{54}627}
\l{}
\l{}
\l{\sou{zenito}. I. (\sou{astr.}) La punto ube la vertikalo di loko iras ren-}
\l{\cel{3}kontre di la sfero cielala. - II. (\sou{metaf.}) La kulmino-punto}
\l{\cel{3}di kozo. - DEFIRS.}
\l{}
\l{\sou{zeolito}. (\sou{mineral.}) Silikato aluminoza hidrata. - DEF.}
\l{}
\l{\sou{zero}.\cel{2}(\sou{matem.}) Cifro, 0-forma, qua - sen valorar ipse - loki-}
\l{\cel{3}zite dextre di altra cifri, igas dekoplan la valoro di ici, per}
\l{\cel{3}suplear la unaji decimala qui mankas. - II. (\sou{fiziko)}. En la}
\l{\cel{3}instrumenti uzata por mezurar la temperaturi, la preso da}
\l{\cel{3}aero, e c., la departo-punto di la grado-konto. - EFIS.}
\l{}
\l{\sou{zesto}. (\sou{bot.)}\cel{2}Parieto membrana qua dividas quarime la internajo}
\l{\cel{3}di nuco. - II. Shelo extera di oranjo e di citrono. - F.}
\l{}
\l{\sou{zezear}. (\sou{netrans.}) Pro defekto, pronuncar le j quale le z.}
\l{\cel{3}e sh quale s. - FS.}
\l{}
\l{\sou{zibelino}. I. (\sou{zool.})\cel{2}Martro di Siberia, kun pili tenua. L.\cel{1}\sou{mar}-}
\l{\cel{3}\sou{tas zibellina}. - II. La felo di ta animalo, uzata kom furo.}
\l{\cel{3}- FIS.}
\l{}
\l{\sou{zigomato}. (\sou{anat.}) Osto di la saliajo di la vango, sub la angulo}
\l{\cel{3}extera di la okulo. - EFIS.}
\l{}
\l{\sou{zigomorfa}. (\sou{bot.)}\cel{3}Dicesas pri la flori qui havas plano di}
\l{\cel{3}simetreso (kom ex.: la violi, la muzelflori).- DEFIS.}
\l{}
\l{\sou{zigosporo}. (\sou{bot.}) Organo genitala sexuoza qua rezultas de la}
\l{\cel{3}fuzo di du gameti izogama o heterogama. - DEFIS.}
\l{}
\l{\sou{zigzago}. Lineo ruptita, qua formacas anguli alterne salianta}
\l{\cel{3}e konkava. - DEFIRS.}
\l{}
\l{\sou{zimazo}. (\sou{biol.})\cel{2}Substanco fermenta, azotoza, nesulfuroza,}
\l{\cel{3}pasable parenta ye la albuminoidi, egardenda kom la reziduo}
\l{\cel{3}di funcionado.- DEFIS.}
\l{}
\l{\sou{zimogeno}. (\sou{biol.}) Substanco quan sekrecas la celuli glandala e}
\l{\cel{3}qua, mixite kun la aquo di la vakuoli, genitas fermento.-}
\l{}
\l{\sou{zimologio}. (\sou{biol.}) La parto di kemio qua traktas fermentaco.-}
\l{\cel{3}DEFIRS.}
\l{}
\l{\sou{zimosimetro}. (\sou{fiziko)}.\cel{2}Instrumento uzata pri mezurar la grado}
\l{\cel{3}di fermentaco di liquido. - DEFIRS.}
\l{}
\l{\sou{zimotekniko.}\cel{1}La arto genitar e direktar la fermentaco. - DEFIRS.}
\l{}
\l{zinko. (\sou{kemio).}\cel{1}Korpo simpla, metala, blue-blanka, kompozajo}
\l{\cel{3}lameloza, duktila, fuzebla, e qua, pos lamino, uzesas por}
\l{\cel{3}facar kovrili di tekti, pluvo-kanali tektala, tubi, balnokuvi,}
\l{\cel{3}siteli, e c. -\cel{1}\sou{Simbolo kemiala}\cel{1}: Zn.}
}
